\documentclass[aps,prl,twocolumn,nofootinbib,floatfix]{revtex4-2}

\usepackage{amsmath,amssymb,amsfonts}
\usepackage{graphicx}
\usepackage{hyperref}
\usepackage{cleveref}

\title{The Source Code of the Universe: Formal Deduction from Discrete Information to Cosmic Density}

\author{Zhang Jun}
\affiliation{New Beginning Universe Research Group (Independent Researcher)}
\email{cnjun939@163.com}
\orcid{0009-0004-9803-3237}

\begin{document}

\maketitle

\begin{abstract}
We present CSQIT (Causal Structure Quantum Information Theory), a fully formalized discrete causal-information axiomatic framework implemented in Lean 4. Starting from a small set of information-theoretic axioms (AxiomA-K), we derive through machine-verifiable proofs a complete deductive chain to a characteristic constant $\theta = 1/(2+2\cos(2\pi/7)) \approx 0.308$, which corresponds to the observed total matter density $\Omega_m = 0.311$ (Planck 2018) with a deviation of approximately 1\%. Among primes $p=3,5,7,11,\ldots$, only $p=7$ yields $\theta(p)$ within the structure formation window $(0.28,0.33)$, establishing a uniqueness constraint from algebraic structure to cosmological parameters. The complete proof code (50 Lean files, ~24,300 lines, 2196 compilation tasks with 0 errors) is available at \url{https://github.com/New-Beginning-Universe-Research-Group/CSQIT}.

\textit{Key results}: (1) Zero-parameter deductive chain from axioms to observable; (2) $\theta \approx \Omega_m$ within $\sim 1\%$; (3) Algebraic uniqueness of $p=7$; (4) Machine-verified formalization; (5) Existential inversion: structure$=7$ as observer existence condition.
\end{abstract}

\section{Introduction}

The reconciliation of quantum mechanics and general relativity remains unresolved. Traditional approaches (string theory, loop quantum gravity, causal set theory) face fundamental challenges: no unique vacuum, no complete semiclassical limit, or no observational anchor. We introduce CSQIT, a fully formalized discrete causal-information framework in Lean 4 that derives a characteristic constant from axioms with zero free parameters, matching cosmological observations within $\sim 1\%$.

\section{Axiomatic System and Formalization}

The CSQIT framework is built upon 10 core axioms (AxiomA-K) formalized in Lean 4:

\begin{itemize}
    \item \textbf{AxiomA}: Weaving structure of rules and relation elements
    \item \textbf{AxiomB}: Causal partial order
    \item \textbf{AxiomC}: Complex-valued quantum amplitude (unitary)
    \item \textbf{AxiomD}: Operational weaving
    \item \textbf{AxiomE-K}: Information capacity, continuum limit, gauge group, etc.
\end{itemize}

All axioms and proofs are implemented and verified in Lean 4 (50 files, ~24,300 lines, 2196 compilation jobs with 0 errors).

\section{Duality Two-One Theorem}

The central theorem states that causal non-triviality and information non-triviality cannot coexist:

\begin{theorem}[Duality Two-One]
In the standard framework, if the output function is non-degenerate (causal non-triviality), then the amplitude function is degenerate (information degeneracy), and vice versa.
\end{theorem}

This establishes a discrete complementarity principle---causal structure and quantum information are two projections of the same underlying reality.

\section{Characteristic Constant $\theta$}

In $\mathbb{Z}_7$ (the smallest prime-order cyclic group satisfying both unitarity and injectivity), the characteristic constant is:

\[
\theta = \frac{1}{2 + 2\cos(2\pi/7)} \approx 0.308
\]

This satisfies the cubic equation $\theta^3 - 6\theta^2 + 5\theta - 1 = 0$. Under the duality interpretation, $\theta$ corresponds to the observed total matter density $\Omega_m = 0.311$ (Planck 2018) with a deviation of approximately 1\%.

\section{Uniqueness of $p=7$}

The algebraic expansion spectrum $\theta(p) = 1/(2+2\cos(2\pi/p))$ is strictly monotonically decreasing for primes $p \ge 3$:

\begin{table}[h]
\centering
\begin{tabular}{ccc}
\hline
$p$ & $\theta(p)$ & In $(0.28,0.33)$? \\
\hline
3 & 1.000 & No \\
5 & 0.382 & No \\
\textbf{7} & \textbf{0.308} & \textbf{Yes} \\
11 & 0.272 & No \\
$\infty$ & 0.250 & No \\
\hline
\end{tabular}
\caption{Algebraic expansion spectrum. Only $p=7$ yields $\theta(p)$ within the structure formation window.}
\label{tab:spectrum}
\end{table}

Structure formation requires $\Omega_m \in (0.28,0.33)$: too low and perturbations cannot collapse into galaxies; too high and the universe recollapses. Only $p=7$ satisfies this constraint.

\section{Existential Inversion}

We reframe the question ``why 7?'' as:

\begin{quote}
``Structure$=7$ is the necessary condition for observers who can ask `why 7' to exist.''
\end{quote}

Let $S(p)$ denote ``a causal lattice with algebraic base $p$ can support internal observers.'' Then $S(p) \implies p=7$. For $p \neq 7$, no observers can exist to question why $p$ is what it is.

\section{Discussion}

CSQIT demonstrates that a formalized axiomatic approach---with zero free parameters---can produce quantitative results comparable to observational data. The proximity of $\theta \approx 0.308$ to $\Omega_m = 0.311$ is the first observational anchor of this structural deduction.

The complete source code is available at \url{https://github.com/New-Beginning-Universe-Research-Group/CSQIT}.

\acknowledgments

This work was conducted independently without external funding.

\bibliography{references}

\end{document}
